\section{Boundary energy profiles and nonlinear parity compatibility}
\label{sec:v9-compatibility}
Fix one of the general channels of Theorem~\ref{thm:v8-main-relative}
and retain observations from both orientations. For an even number of
flights, the first and last impacts have the same contact type; for an
odd number they have different types. Write
$\mathcal F_{00},\mathcal F_{11}$ for the two even limiting functions
and $\mathcal F_{01}=\mathcal F_{10}$ for the odd one. Their
normalization is that of \eqref{eq:v8-relative-integral}, with the
same free area $A$ and multiplier $\gamma$. Thus they are scalar
functions derived from selected physical event probabilities, not
postulated independent boundary densities.

\subsection{An intrinsic positive factorization}
For each contact type put $a_b=S_b''(0)>0$ and introduce the signed
Morse coordinate and its increasing inverse,
\[
 t_b(u)=\operatorname{sgn}(u)\sqrt{2S_b(u)},\qquad u=\chi_b(t).
\]
The value at $u=0$ is defined by smooth extension. Set
\begin{equation}\label{eq:v9-energy-profile}
 \begin{split}
 f_b(t)&=B_b(\chi_b(t))\chi_b'(t),\\
 \mathcal V_b(x)&=\frac{\sqrt{a_b}}2
            \{f_b(\sqrt{2x})+f_b(-\sqrt{2x})\},\qquad x\ge0.
 \end{split}
\end{equation}
We call $\mathcal V_b$ the normalized boundary energy profile. It is
a symmetrized weighted energy datum: no assertion that it determines
a nonsymmetric contact graph is implicit in this terminology.

\begin{theorem}[Nonlinear boundary compatibility]
\label{thm:v9-compatibility}
On a sufficiently small common closed collar the functions
$\mathcal V_b$ are smooth, positive, and satisfy $\mathcal V_b(0)=1$.
For $b,c\in\{0,1\}$ the limiting physical laws have the representation
\begin{equation}\label{eq:v9-energy-integral}
 \mathcal F_{bc}(d)=\frac{2}{\pi d^2}
 \int_{\substack{x,y>0\\x+y<d}}
       (d-x-y)\frac{\mathcal V_b(x)\mathcal V_c(y)}{\sqrt{xy}}
                                               \,\dd x\dd y.
\end{equation}
Let $*$ denote Volterra convolution on the positive half-line, used
only at arguments in this collar, and define
\begin{equation}\label{eq:v9-kernel-def}
 k_b(x)=x^{-1/2}\mathcal V_b(x),\qquad
 K_{bc}(d)=\frac\pi2\frac{\dd^2}{\dd d^2}
                                \{d^2\mathcal F_{bc}(d)\}.
\end{equation}
Then $K_{bc}=k_b*k_c$, $K_{bc}(0)=\pi$, and
\begin{equation}\label{eq:v9-volterra-compatibility}
                  K_{01}*K_{01}=K_{00}*K_{11}.
\end{equation}
Within the class of smooth normalized profiles, $\mathcal F_{00}$
and $\mathcal F_{11}$ determine $\mathcal V_0,\mathcal V_1$
uniquely and therefore determine the complete odd law
$\mathcal F_{01}$ on the same collar. The assertions hold without
equal-curvature, even-graph or symmetry assumptions. All the forward
identities have the mixed geometric regularity of the relative law.
\end{theorem}
\begin{proof}
Strict convexity makes the signed Morse coordinate smooth and
increasing near zero, with derivative $\sqrt{a_b}$ at zero.
Consequently $f_b$ is smooth, positive and has value $a_b^{-1/2}$.
Its even part is a smooth function of $t^2$, by the even-function
argument in Proposition~\ref{prop:v9-morse-transport}. This proves
the asserted properties of $\mathcal V_b$, including its parameter
regularity. Reversing the choice of transverse direction merely
interchanges the two terms in \eqref{eq:v9-energy-profile}.

Under $x=S_b(u)=t^2/2$, the pushforward of $B_b(u)\,\dd u$ has,
for $x>0$, density
\[
 \frac{f_b(\sqrt{2x})+f_b(-\sqrt{2x})}{\sqrt{2x}}
       =\sqrt{\frac2{a_b}}\,x^{-1/2}\mathcal V_b(x).
\]
Insert this expression twice into
\eqref{eq:v8-boundary-integral}, whose prefactor for types $b,c$
is $\sqrt{a_ba_c}/(\pi d^2)$. The Hessian factors cancel and give
\eqref{eq:v9-energy-integral}. This calculation also proves equality
of the two odd orientations. It is valid for each even orientation
by applying the relative law after interchanging the contact types.

The kernels $k_b$ are locally integrable. Fubini's theorem rewrites
\eqref{eq:v9-energy-integral} as
\[
 \frac\pi2d^2\mathcal F_{bc}(d)
                  =\int_0^d(d-s)(k_b*k_c)(s)\,\dd s.
\]
The convolution is continuous down to zero: substitution $x=ds$
gives
\[
 (k_b*k_c)(d)=\int_0^1
       \frac{\mathcal V_b(ds)\mathcal V_c(d(1-s))}{\sqrt{s(1-s)}}
                                                         \,\dd s,
\]
whose value at zero is $\pi$. Differentiation proves
\eqref{eq:v9-kernel-def}. Associativity and commutativity of these
locally integrable convolutions, again by Fubini on a bounded
simplex, give \eqref{eq:v9-volterra-compatibility}. No extension of
the physical law beyond the collar is used.

It remains to prove uniqueness, since a formal square root of Taylor
coefficients alone would not suffice for smooth functions. Suppose
$k=x^{-1/2}V(x)$ and $\ell=x^{-1/2}\widetilde V(x)$, with smooth
$V(0)=\widetilde V(0)=1$, have $k*k=\ell*\ell$ on the collar.
Put $f=k-\ell$, $a=k+\ell$ and $k_*(x)=x^{-1/2}$. Then
$a*f=0$ and
\[
 a(x)=2x^{-1/2}+x^{1/2}w(x)
\]
for a smooth $w$. The identity $k_* * k_* =\pi$ gives
\[
 (k_* * a)(x)=2\pi+R(x),\qquad
 R(x)=x\int_0^1\frac{s^{1/2}}{(1-s)^{1/2}}w(xs)\,\dd s.
\]
Here $R(0)=0$ and $R'$ is bounded. Convolving $a*f=0$ with
$k_*$ and differentiating its absolutely continuous integral yields
\[
                    2\pi f+R'*f=0
\]
almost everywhere. Hence
$|f(x)|\le C\int_0^x|f(s)|\,\dd s$.
The absolutely continuous function $\int_0^x|f(s)|\,\dd s$ has
zero initial value, and the elementary Gronwall inequality makes it
zero throughout the collar. Thus $k=\ell$. Applying this argument
to each of $K_{00}$ and $K_{11}$ proves uniqueness of both profiles;
\eqref{eq:v9-energy-integral} then determines the odd law.
\end{proof}

\begin{theorem}[Stability of the complete boundary energy profile]
\label{thm:v9-profile-stability}
Fix $D>0$ in the common collar. Let $V,\widetilde V$ be two smooth
normalized profiles, with $V(0)=\widetilde V(0)=1$ and
$\|V'\|_\infty,\|\widetilde V'\|_\infty\le B$ on $[0,D]$.
Let $F,\widetilde F$ be their same-type laws defined by
\eqref{eq:v9-energy-integral}, and put
$K=(\pi/2)(d^2F)''$, $\widetilde K=(\pi/2)(d^2\widetilde F)''$.
Then
\begin{equation}\label{eq:v9-profile-stability}
 \|V-\widetilde V\|_{C^0([0,D])}
 \le \frac D\pi e^{BD/2}\|K'-\widetilde K'\|_{C^0([0,D])}
 \le C_{B,D}\|F-\widetilde F\|_{C^3([0,D])}.
\end{equation}
Thus the two complete boundary profiles are locally Lipschitz
functions of the two same-type laws on their physical image, when
the latter are given this specified differentiated norm. The resulting
odd law is Lipschitz in $C^0$ on uniformly bounded profile sets.
\end{theorem}
\begin{proof}
Use $k=x^{-1/2}V$, $\ell=x^{-1/2}\widetilde V$,
$f=k-\ell$, $a=k+\ell$, $k_*=x^{-1/2}$, and put
$Q=K-\widetilde K=a*f$. The convolution in the preceding proof
has the more direct expression
\[
 (k_* * a)(x)=2\pi+R(x),\qquad
 R(x)=\int_0^1\frac{V(xs)+\widetilde V(xs)-2}{\sqrt{s(1-s)}}
                                                        \,\dd s.
\]
It follows that $R(0)=0$ and
\[
 \|R'\|_\infty\le
  2B\int_0^1\frac{s}{\sqrt{s(1-s)}}\,\dd s=\pi B.
\]
Since $Q(0)=0$, differentiation under a convolution, or integration
by parts before differentiation, gives
\[
 2\pi f+R'*f=(k_* * Q)'=k_* * Q'.
\]
The right side is bounded by $2\sqrt{x}\,\|Q'\|_\infty$.
The normalized initial values imply that $f$ extends continuously
with $f(0)=0$. Therefore
\[
 |f(x)|\le\frac{\sqrt D}{\pi}\|Q'\|_\infty
                          +\frac B2\int_0^x|f(s)|\,\dd s.
\]
Gronwall's inequality gives
$\|f\|_\infty\le\sqrt D\,e^{BD/2}\|Q'\|_\infty/\pi$.
Multiplication by $\sqrt x$ proves the first estimate. For the second,
\[
 Q'(d)=\frac\pi2\{6(F-\widetilde F)'(d)
       +6d(F-\widetilde F)''(d)+d^2(F-\widetilde F)'''(d)\},
\]
which is bounded by the asserted $C^3$ norm. Finally the integral
\eqref{eq:v9-energy-integral} is a bilinear map of the profiles with
a positive kernel of normalized total mass one. Subtracting its two
products proves the $C^0$ Lipschitz assertion for the odd law.
\end{proof}

This is stability of normalized energy profiles, not of an arbitrary
nonsymmetric contact graph. It is an inverse on the image of the
forward profile map; no claim that every positive scalar function is
a convolution square is needed. The derivative norm in
\eqref{eq:v9-profile-stability} is part of the observation statement.
The unregularized theorem does not replace it by uniform error of raw
probabilities. The finite-preparation theorem in
Section~\ref{sec:v10-acquisition} adds regularization and supplied
smoothness bounds to address that different data problem. For finite
jets, an explicit polynomial
inverse is available without solving a function-valued equation.

\subsection{Finite jets and observable identities}
For a smooth right germ $h$, write $[d^n]h=h^{(n)}(0)/n!$.
Fix $M\ge0$ and work in the truncated polynomial ring
$\R[z]/(z^{M+1})$. Introduce
\begin{equation}\label{eq:v9-jet-transform}
 \begin{split}
 \mathcal M_{bc}^{[M]}(z)
   &=\sum_{n=0}^M\frac{(n+2)!}{2}[d^n]\mathcal F_{bc}\,z^n,\\
 \mathcal A_b^{[M]}(z)
   &=\sum_{n=0}^M (\tfrac12)_n[x^n]\mathcal V_b\,z^n,
 \qquad (\tfrac12)_0=1.
 \end{split}
\end{equation}
The notation $(1/2)_n$ denotes the rising factorial. Truncation is
important: even for analytic input, no convergence radius of a
factorially reweighted infinite series is being asserted.

\begin{theorem}[Rank-one jets and stable finite-order inversion]
\label{thm:v9-jet-inverse}
The transformed jets satisfy
\begin{equation}\label{eq:v9-rank-one}
 \mathcal M_{bc}^{[M]}=\mathcal A_b^{[M]}\mathcal A_c^{[M]},
 \qquad
 (\mathcal M_{01}^{[M]})^2
                   =\mathcal M_{00}^{[M]}\mathcal M_{11}^{[M]}
                   \pmod {z^{M+1}}.
\end{equation}
Both profiles' order-$M$ jets are recovered from the two same-type
laws by the following triangular recursion. If
$m_{bb,n}=[z^n]\mathcal M_{bb}^{[M]}$ and
$\alpha_{b,n}=[z^n]\mathcal A_b^{[M]}$, then
\begin{equation}\label{eq:v9-square-root-recursion}
 \alpha_{b,0}=1,\qquad
 \alpha_{b,n}=\frac12\left(m_{bb,n}
       -\sum_{r=1}^{n-1}\alpha_{b,r}\alpha_{b,n-r}\right),
 \qquad [x^n]\mathcal V_b=\frac{\alpha_{b,n}}{(1/2)_n}.
\end{equation}
For fixed $M$ these reconstruction maps, and the resulting prediction
of the odd jet, are Lipschitz on bounded sets of input coefficients.
They are smooth with respect to geometric parameters on the compact
families of the relative law.
\end{theorem}
\begin{proof}
Write $v_{b,r}=[x^r]\mathcal V_b$. The beta integral on the
simplex gives, for $r,s\ge0$,
\[
 \int_{x,y>0,\ x+y<d}(d-x-y)x^{r-1/2}y^{s-1/2}\,\dd x\dd y
       =\frac{\pi(1/2)_r(1/2)_s}{(r+s+2)!}\,d^{r+s+2}.
\]
This follows by first setting $x=dX,y=dY$ and integrating in $Y$,
then in $X$; both integrals are the ordinary beta integral.
Taylor's integral remainders, bounded on the scaled simplex, justify
coefficient extraction for smooth rather than only analytic profiles.
Thus
\begin{equation}\label{eq:v9-beta-coefficients}
 [d^n]\mathcal F_{bc}
       =\frac2{(n+2)!}\sum_{r+s=n}
                    (1/2)_r(1/2)_s v_{b,r}v_{c,s}.
\end{equation}
Multiplying by $(n+2)!/2$ gives the first identity in
\eqref{eq:v9-rank-one}; its determinant identity is an immediate
polynomial consequence. All constant coefficients are one. Comparing
the degree-$n$ coefficient of
$\mathcal M_{bb}^{[M]}=(\mathcal A_b^{[M]})^2$ leaves
$2\alpha_{b,n}$ as the only unknown term, proving
\eqref{eq:v9-square-root-recursion}. Induction makes every output
coefficient a polynomial in finitely many input coefficients, divided
only by fixed positive numbers. Its derivative is bounded on each
bounded coefficient box, proving the Lipschitz assertion. Finite
parameter differentiation of these polynomial identities proves the
last assertion without any infinite-order uniformity claim.
\end{proof}

\begin{corollary}[The first two nonlinear compatibility identities]
\label{cor:v9-first-two}
With primes denoting right offset derivatives at zero,
\begin{align}
 2\mathcal F_{01}'&=\mathcal F_{00}'+\mathcal F_{11}',
                                  \label{eq:v9-first-identity}\\
 \mathcal F_{01}''&=
       \frac{\mathcal F_{00}''+\mathcal F_{11}''}{2}
       -\frac3{16}(\mathcal F_{00}'-\mathcal F_{11}')^2.
                                  \label{eq:v9-second-identity}
\end{align}
In particular the odd second derivative is at most the mean of the
two even second derivatives, with the exact deficit displayed above.
\end{corollary}
\begin{proof}
The first two nonconstant coefficients of $\mathcal M_{bc}$ are
$3\mathcal F_{bc}'$ and $6\mathcal F_{bc}''$.
Equate degrees one and two in
$\mathcal M_{01}^2=\mathcal M_{00}\mathcal M_{11}$ and substitute
the degree-one equality in the degree-two equality. This gives the
two formulas, including the negative square coefficient.
\end{proof}

These identities concern nonlinear offset coefficients. The
zeroth-order condition merely says that every normalized law has
value one. Neither the zeroth-order condition nor the multiplier
alone determines the first two coefficients. The geometric variation
in Proposition~\ref{prop:v4-quartic} and the realization in
Theorem~\ref{thm:v4-jet-fiber} exhibit nonconstant boundary data at
fixed leading selected-channel invariants. The compatibility theorem
organizes such data without imposing the identical-even hypothesis
of the separate graph-jet inverse in the appendix.

\subsection{Finite physical records and a quantitative defect}
The preceding identities are exact statements about the limiting
functions. The relative law also controls how they can be read from
long but finite records. For $n\ge1$ define
\begin{equation}\label{eq:v9-finite-laws}
 \begin{split}
 F_{bb,n}(d)&=\frac{2A\sinh(2n\gamma)}{d^2}
              \Pr(E_{2n,e_b}(d)),\quad e_0=e,\ e_1=\bar e,\\
 F_{01,n}(d)&=\frac{2A\sinh((2n+1)\gamma)}{d^2}
              \Pr(E_{2n+1,e}(d)).
 \end{split}
\end{equation}
They are extended to $d=0$ by their exact smooth right limits.
These definitions use the selected patch labels and the normalization
$A,\gamma$; they are not unlabelled binary data with unknown
normalizer. Estimation of these auxiliary quantities is a separate
observation problem.

\begin{corollary}[Finite-record compatibility and profile-jet recovery]
\label{cor:v9-finite-records}
Let $M$ and the number of geometric derivatives be fixed. Apply the
transform \eqref{eq:v9-jet-transform} to the three finite laws in
\eqref{eq:v9-finite-laws}, obtaining $\mathcal M_{bc,n}^{[M]}$.
With the maximum coefficient norm on truncated polynomials,
\begin{equation}\label{eq:v9-finite-defect}
 \left\|(\mathcal M_{01,n}^{[M]})^2
         -\mathcal M_{00,n}^{[M]}\mathcal M_{11,n}^{[M]}\right\|
                                       \le C_M\tau^{2n}.
\end{equation}
The same estimate holds for each fixed mixed geometric derivative.
Suppose the measured order-$M$ jets of the two same-type finite laws
have coefficient error at most $\epsilon$, with their known constant
coefficients fixed to one. The recursion
\eqref{eq:v9-square-root-recursion} reconstructs both boundary profile
jets, and predicts the odd limiting jet, with error
\begin{equation}\label{eq:v9-finite-recovery}
                           C_M(\epsilon+\tau^{2n}).
\end{equation}
For every fixed subcollar $[0,D]\subset[0,d_*]$, the finite kernels
$K_{bc,n}=(\pi/2)(d^2F_{bc,n})''$ also satisfy
\[
 \|K_{01,n}*K_{01,n}-K_{00,n}*K_{11,n}\|_{C^0([0,D])}
                                     \le C_D\tau^{2n}.
\]
\end{corollary}
\begin{proof}
Theorem~\ref{thm:v8-main-relative}, applied to each orientation and
parity, bounds $F_{bc,n}-\mathcal F_{bc}$ in every fixed $C^k$
norm by $C_k\tau^{2n}$. Coefficient extraction and the finite
transform preserve this order. The determinant in
\eqref{eq:v9-rank-one} is a polynomial in bounded coefficients;
subtract its value at the limiting jets to prove
\eqref{eq:v9-finite-defect}. The product rule proves the parameter
derivative assertion. Apply the finite polynomial inverse of
Theorem~\ref{thm:v9-jet-inverse} to the two measured jets to get
\eqref{eq:v9-finite-recovery}; this is also defined for noisy jets
outside the physical image. Finally the relative $C^2$ estimate
bounds $K_{bc,n}-K_{bc}$ uniformly on $[0,D]$. For bounded
functions there, $\|f*g\|_\infty\le D\|f\|_\infty\|g\|_\infty$.
Telescoping the four products and using
\eqref{eq:v9-volterra-compatibility} proves the last claim.
\end{proof}

The two oriented even laws suffice to reconstruct both normalized
boundary energy profiles and to predict the odd law. Odd records are
needed for the independent compatibility test, not as a third datum
for this reconstruction. Neither conclusion recovers full collision
arrays or an arbitrary billiard table. Unlike the fixed-flight
parametric inverse below, the compatibility defect and its vanishing
rate use the long-bridge relative theorem directly.
Section~\ref{sec:v10-acquisition} supplies a regularized finite-preparation
version of the full-profile inverse, with the required normalization
and smoothness bounds stated explicitly.
